\section{Explicit Lifetimes and Fixed Fallbacks}
\label{app:lifecycle}
Let claim $c$ start at $x_c\in[\ell_c,u_c]$ and end at $z_c\in[v_c,w_c]$.
Assume $u_c<v_c$ and independent start and end dates.
Its explicit lifetime is $[x_c,z_c)$, additionally shortened by accepted direct replacements.
Introduce a hidden event $e_c$ with bounds $[v_c,w_c]$ and the pair $e_c\to c$.
Every world has $x_c<z_c$.
The hidden event retires $c$ exactly when the explicit end occurs.
It therefore preserves activity in every world.
The support boundaries update to
\[
 p'_c=\min(p_c,w_c),\qquad r'_c=\min(r_c,v_c).
\]
The same support formulas apply, with at most two witness records per claim.
Hidden end events never enter a reader's context.

The strict separation condition is necessary for this independent-date reduction.
For start bounds $[0,2]$ and end bounds $[1,3]$, the product permits start 2 and end 1.
A valid lifetime excludes that assignment, but direct replacement ignores an earlier event.
The adapter therefore leaves overlapping start--end bounds unmodeled.
It also retains approximate dates without justified finite bounds as unmodeled text.

Let $F$ contain fixed fallback items and $S_x$ contain the modeled support set.
The operational mask is $S_x\cup F$.
Its union over worlds is $P\cup F$, and its intersection is $H\cup F$.
The context test therefore compares $\operatorname{Top}_k(P\cup F)$ with $\operatorname{Top}_k(H\cup F)$.
Fallback eligibility records a policy choice and never receives the guaranteed-support tag.

Additional source constraints can restrict the independent date rectangle.
Possible support then remains an outer bound, while guaranteed support remains an inner bound.
Equal endpoint contexts still certify stability under the extra constraints.
Different endpoint contexts need not establish instability within that smaller family.
Counterworld exactness uses the stated independent-date model.
